\documentclass[11pt]{article}
\usepackage[english]{babel}
\usepackage[a4paper,margin=0.9in]{geometry}
\usepackage[T1]{fontenc}
\usepackage[utf8]{inputenc}
\usepackage{lmodern}
\usepackage{amsmath,amssymb,mathtools,mathrsfs}
\usepackage{microtype}
\usepackage{fancyhdr}
\usepackage{array,longtable,enumitem,multicol}
\usepackage{tikz}
\setlength{\parindent}{1.5em}
\setlength{\parskip}{0.18em}
\emergencystretch=6em
\pagestyle{fancy}
\fancyhf{}
\cfoot{\thepage}
\newcommand{\Jac}[2]{\left(\frac{#1}{#2}\right)}
\newcommand{\disc}{\operatorname{disc}}
\newcommand{\Norm}{\operatorname{N}}
\providecommand{\dd}{\,d}
\providecommand{\tht}{\vartheta}
\providecommand{\vark}{\varkappa}
\providecommand{\sn}{\operatorname{sn}}
\providecommand{\cn}{\operatorname{cn}}
\providecommand{\dn}{\operatorname{dn}}
\providecommand{\am}{\operatorname{am}}
\providecommand{\ctg}{\operatorname{cotg}}
\providecommand{\cotg}{\operatorname{cotg}}
\providecommand{\eps}{\varepsilon}
\providecommand{\ii}{\mathrm{i}}
\providecommand{\tg}{\operatorname{tg}}
\providecommand{\jac}[2]{\left(\frac{#1}{#2}\right)}
\newcommand{\Om}{\Omega}
\newcommand{\Omm}{\Omega_m}
\newcommand{\Hm}{H_m}
\newcommand{\Lm}{\Omega_m}
\newcommand{\Normone}{\mathrm{N}_1}
\newcommand{\frp}{\mathfrak p}
\newcommand{\frP}{\mathfrak P}
\newcommand{\frG}{\mathfrak G}
\newcommand{\frA}{\mathfrak A}
\newcommand{\frB}{\mathfrak B}
\newcommand{\Gg}{\mathfrak G}
\newcommand{\Aa}{\mathfrak A}
\newcommand{\Bb}{\mathfrak B}
\newcommand{\pp}{\mathfrak p}
\newcommand{\G}{\mathfrak G}
\newcommand{\B}{\mathfrak B}
\newcommand{\A}{\mathfrak A}
\newcommand{\frakA}{\mathfrak a}
\newcommand{\Na}{N_a}
\newcommand{\abs}[1]{\left|#1\right|}
\newcommand{\idxitem}[1]{\par\hangindent=1.4em\hangafter=1 #1}
\providecommand{\Lzero}{\mathfrak L_0}
\providecommand{\Kgrp}{\mathfrak K}
\lhead{Weber Vol. III §§87-89}
\rhead{English translation}
\begin{document}

\section*{Tenth Section. Algebraic numbers and forms.}

\section*{§ 87. Ideals and forms in algebraic fields.}

The interest which the theory of elliptic functions has for the algebraist arises from its application to the theory of quadratic irrational numbers, to which it stands in an analogous relation as the roots of unity stand to the rational numbers. It offers us the first, and until now the only, example, more exactly known in its laws, of a domain of algebraic numbers which goes beyond the cyclotomic numbers. In order to present the theory of these numbers more closely, we must first place before it some theorems from the theory of forms and algebraic fields, and especially of quadratic forms.

In Vol. II, §§ 163, 169, a connection was obtained between the ideals of an algebraic field $\Omega$ of degree $n$ and certain homogeneous functions of degree $n$ in $n$ variables. Before applying this to quadratic fields, let us briefly recall the general theorems.

Let $\mathfrak a$ be an ideal of a field $\Omega$, and let $\alpha_1,\alpha_2,
\ldots,\alpha_n$ be a basis of this ideal. Further let $\omega_1,\omega_2,
\ldots,\omega_n$ be a minimal basis of $\Omega$. Then
\begin{equation}
 (\alpha_1,\alpha_2,\ldots,\alpha_n)=(A)(\omega_1,\omega_2,
\ldots,\omega_n),\tag{1}
\end{equation}
where $(A)$ denotes a linear substitution with integral coefficients. To pass to another basis $\alpha'_1,\alpha'_2,
\ldots,\alpha'_n$ of $\mathfrak a$, make a linear substitution $L$ with determinant $\pm1$,
\begin{equation}
 (\alpha'_1,\alpha'_2,
\ldots,\alpha'_n)=(L)(\alpha_1,\alpha_2,
\ldots,\alpha_n),\tag{2}
\end{equation}
whence
\begin{equation}
 (\alpha'_1,\alpha'_2,
\ldots,\alpha'_n)=(L)(A)(\omega_1,\omega_2,
\ldots,\omega_n).\tag{3}
\end{equation}
By the norm $N(\mathfrak a)$ of the ideal $\mathfrak a$ one understands the absolute value of the determinant $A$ of the substitution $(A)$:
\[
 A=\sum\pm a_{1,1}a_{2,2}\cdots a_{n,n},
\]
and by (3) the norm is independent of the choice of basis.

1. We shall call the basis $(\alpha_\nu)$ positive if the determinant $A$ is positive, and therefore equal to the norm of $\mathfrak a$. If $(\alpha_\nu)$ is positive, then $(\alpha'_\nu)$ is positive precisely when the determinant of the substitution is $L=+1$.

From every basis one can derive a positive basis by a substitution with determinant $-1$, for example by interchanging two of the $\alpha_\nu$, and in what follows we shall mostly use only positive bases.

The basis determined in Vol. II, § 163, (3), is positive. Let $t_1,t_2,
\ldots,t_n$ be a system of independent variables and
\begin{equation}
 \lambda=\sum_r\alpha_rt_r\tag{4}
\end{equation}
a basis form with positive basis. Then (Vol. II, § 164)
\begin{equation}
 N(\lambda)=N(\mathfrak a)T,\tag{5}
\end{equation}
where
\begin{equation}
 T=\sum\pm t_{1,1}t_{2,2}\cdots t_{n,n}\tag{6}
\end{equation}
is a primitive integral homogeneous form of degree $n$ in the variables $t_\nu$. The $t_{r,s}$ are linear functions of the $t_\nu$.

Applying the substitution (2) to (4), one obtains
\begin{equation}
 \lambda=\lambda'=\sum_r\alpha'_rt'_r,\tag{7}
\end{equation}
where
\begin{equation}
 (t_1,t_2,\ldots,t_n)=L'(t'_1,t'_2,
\ldots,t'_n),\tag{8}
\end{equation}
and $L'$ is the transposed substitution of $L$. If $(\alpha'_\nu)$ is likewise a positive basis, then $L'$ has determinant $+1$, and formula (5) gives
\begin{equation}
 N(\lambda)=N(\mathfrak a)T',\tag{9}
\end{equation}
where $T'$ is a homogeneous function of degree $n$ in the variables $t'$ which arises from $T$ by the substitution (8).

Two forms which arise from one another by an integral linear substitution with determinant $+1$ are called equivalent. Sometimes the forms are also called equivalent when the determinant of the substitution is $-1$, but then with the qualification ``improperly.'' Thus (8) proves:

2. If we call $T$ the form belonging to the basis $(\alpha_\nu)$ of $\mathfrak a$, then the forms belonging to the different positive bases of the same ideal are equivalent.

Denote the conjugate values of $\alpha_s,\omega_s$ by
\[
 \alpha_{s,1},\alpha_{s,2},\ldots,\alpha_{s,n};\qquad
 \omega_{s,1},\omega_{s,2},\ldots,\omega_{s,n},
\]
and put
\begin{equation}
 \alpha_{s,r}=\sum_\nu a_{s,\nu}\omega_{\nu,r}.\tag{10}
\end{equation}
Then
\begin{equation}
 \left(\sum\pm\alpha_{1,1}\alpha_{2,2}\cdots\alpha_{n,n}\right)^2=N(\mathfrak a)^2\Delta,\tag{11}
\end{equation}
where $\Delta$ is the fundamental number of the field $\Omega$, namely the square of the determinant
\[
 \Delta=\left(\sum\pm\omega_{1,1}\omega_{2,2}\cdots\omega_{n,n}\right)^2
\]
(Vol. II, § 162).

If $\lambda_1,\lambda_2,
\ldots,\lambda_n$ are the conjugate values of $\lambda$, then by (4)
\begin{equation}
 \lambda_r=\sum_\nu\alpha_{\nu,r}t_\nu\tag{12}
\end{equation}
is a linear substitution for the variables $t_\nu$ with substitution determinant
\begin{equation}
 r=N(\mathfrak a)\sqrt{\Delta},\tag{13}
\end{equation}
and through this, by (5), the form $T$ passes into the product
\begin{equation}
 T'=[N(\mathfrak a)]^{-1}\lambda_1\lambda_2\cdots\lambda_n,\tag{14}
\end{equation}
because $N(\lambda)$ is the product of the conjugate values of $\lambda$ (Vol. II, § 151). We can apply invariant theory to this (Vol. I, § 65), and if we form the Hessian determinant
\begin{equation}
 H=\det\left(\frac{\partial^2T}{\partial t_i\partial t_j}\right),\tag{15}
\end{equation}
then
\begin{equation}
 H'=r^{n-2}H.\tag{16}
\end{equation}
If, however, we form $H'$ from (14) in the variables $\lambda_\nu$ and observe that an $n$-rowed determinant whose diagonal terms are $0$ and all other terms are $1$ has the value $(-1)^{n-1}(n-1)$,
\footnote{One can easily prove this theorem by adjoining an $(n+1)$st row and an $(n+1)$st column to the determinant, in such a way that in the adjoined row only ones occur and in the column, except for the last element, zeros occur.}
then
\[
 H'=(-1)^{n-1}(n-1)N(\mathfrak a)^{-n}(\lambda_1\lambda_2\cdots\lambda_n)^{n-2}
 =(-1)^{n-1}(n-1)N(\mathfrak a)^{-2}T^{n-2},
\]
and consequently, by (13) and (16),
\begin{equation}
 H=(-1)^{n-1}(n-1)T^{n-2}\Delta.\tag{17}
\end{equation}

\section*{§ 88. Ideal classes and form classes.}

In Vol. II, § 170, we explained the equivalence of ideals and the ideal classes. Accordingly two ideals $\mathfrak a,\mathfrak b$ were equivalent if there is an integral or fractional number of the field $\Omega$ which is the quotient of $\mathfrak b$ and $\mathfrak a$, that is,
\begin{equation}
 \eta\mathfrak a=\mathfrak b,\tag{1}
\end{equation}
and it was shown that the ideals are thereby divided into classes, and that the number of these classes is finite. In some cases it is expedient to take the concept of equivalence somewhat more narrowly and to call $\mathfrak a$ and $\mathfrak b$ equivalent only when there is a number $\eta$ satisfying condition (1) with positive norm. This distinction naturally does not arise when the norms of all numbers are positive, as for instance in the imaginary quadratic field. It also does not arise when there are units of norm $-1$, since, if $\varepsilon$ is such a unit, then either $N(\eta)$ or $N(\varepsilon\eta)$ is positive, and $\varepsilon$ and $\varepsilon\eta$ simultaneously satisfy condition (1).

But if there are no units whose norm is $-1$, although there are other numbers in $\Omega$ with negative norm, then under the narrower definition of equivalence each ideal class divides once more into two classes. We shall keep here to the narrower concept of equivalence, which for example is in force in some real quadratic fields.
\footnote{Namely, when Pell's equation $t^2-Du^2=-4$ has no solution.}

In (1), $\eta$ may be fractional. But if $\alpha$ is an integral number divisible by $\mathfrak a$, then $\eta\alpha$ is an integral number divisible by $\mathfrak b$, and from this one obtains the following. If
\begin{equation}
 (\alpha_1,\alpha_2,
\ldots,\alpha_n)\tag{2}
\end{equation}
is a basis of $\mathfrak a$, then
\[
 (\beta_1,\beta_2,
\ldots,\beta_n)=(\alpha_1\eta,\alpha_2\eta,
\ldots,\alpha_n\eta)
\]
is a basis of $\mathfrak b$; and if the first basis is positive, then, under the narrower concept of equivalence, so is the second.

For the necessary and sufficient condition that (2) be a basis of $\mathfrak a$ is that every integral number $\alpha$ divisible by $\mathfrak a$, and only such a number, be contained in the form
\begin{equation}
 \alpha=x_1\alpha_1+x_2\alpha_2+\cdots+x_n\alpha_n\tag{3}
\end{equation}
with integral $x_1,x_2,
\ldots,x_n$. Consequently all numbers
\begin{equation}
 \beta=x_1\alpha_1\eta+x_2\alpha_2\eta+\cdots+x_n\alpha_n\eta\tag{4}
\end{equation}
are divisible by $\mathfrak b$. Conversely, if $\beta$ is an integral number divisible by $\mathfrak b$, then $\beta/\eta$ is divisible by $\mathfrak a$ and hence representable in the form (3). Thus $\beta$ is representable by (4).

Furthermore, using the notation of (9) and (10), § 87, when $\alpha$ is a positive basis:
\[
 \sum\pm \alpha_{1,1}\alpha_{2,2}\cdots\alpha_{n,n}=N(\mathfrak a)\sum\pm \omega_{1,1}\omega_{2,2}\cdots\omega_{n,n},
\]
\[
 \sum\pm \beta_{1,1}\beta_{2,2}\cdots\beta_{n,n}=N(\eta)
 \sum\pm \alpha_{1,1}\alpha_{2,2}\cdots\alpha_{n,n}
 =\pm N(\mathfrak b)
 \sum\pm \omega_{1,1}\omega_{2,2}\cdots\omega_{n,n},
\]
and since by (1)
\begin{equation}
 N(\mathfrak b)=N(\eta)N(\mathfrak a),\tag{5}
\end{equation}
the positive sign must stand at $\pm N(\mathfrak b)$.

If $\lambda$ is a basis form of $\mathfrak a$, then $\eta\lambda$ is a basis form of $\mathfrak b$, and
\begin{equation}
 \begin{aligned}
 N(\lambda)&=N(\mathfrak a)T,\\
 N(\eta\lambda)&=N(\mathfrak b)T,
 \end{aligned}\tag{6}
\end{equation}
and hence the same form $T$ belongs to the two ideals $\mathfrak a,\mathfrak b$. If equivalent forms $T$ are combined into a class, we can say by 2.:

3. To every ideal class there belongs a determinate form class.

In answering the question whether different ideal classes can belong to one and the same form class, we make the assumption that $\Omega$ is a normal field. Suppose, then, that equivalent forms $T$ and $T'$ belong to two ideals $\mathfrak a$ and $\mathfrak a'$. If $\lambda$ is a basis form of $\mathfrak a$, then
\begin{equation}
 N(\lambda)=N(\mathfrak a)T,\tag{7}
\end{equation}
and if we carry $T'$ into $T$ by a linear substitution, and understand by $\lambda'$ a basis form of $\mathfrak a'$ (with the same variables $t$ as $\lambda$), then
\begin{equation}
 N(\lambda')=N(\mathfrak a')T.\tag{8}
\end{equation}
The linear factors of $N(\lambda)$ and $N(\lambda')$ must therefore agree with one another, apart from a constant factor, since an integral function of arbitrary variables can be decomposed into irreducible factors (here linear factors) in only one way (Vol. I, § 20, Vol. II, § 152); and since all these factors belong to the same field $\Omega$, there is among the factors conjugate to $\lambda'$ one which satisfies the condition
\[
 \lambda=\eta\lambda',
\]
where $\eta$ is a number in $\Omega$. The functionals $\lambda$ and $\lambda'$, and therefore the corresponding ideals, are thus equivalent.

4. If, in a normal field, the same form class belongs to two ideals $\mathfrak a$ and $\mathfrak a'$, then among the ideals conjugate to $\mathfrak a'$ there is one which is equivalent to $\mathfrak a$.

\section*{§ 89. Composition of forms and multiplication of ideals.}

Let $\mathfrak a$ and $\mathfrak b$ now be two ideals in $\Omega$ and let
\begin{equation}
 \mathfrak c=\mathfrak a\mathfrak b\tag{1}
\end{equation}
be the product formed from them. Further let
\begin{equation}
\begin{aligned}
 \mathfrak a&=(\alpha_1,\alpha_2,
\ldots,\alpha_n),\\
 \mathfrak b&=(\beta_1,\beta_2,
\ldots,\beta_n),\\
 \mathfrak c&=(\gamma_1,\gamma_2,
\ldots,\gamma_n)
\end{aligned}\tag{2}
\end{equation}
be bases of these ideals. Since each number $\alpha_h\beta_k$ belongs to $\mathfrak c$, there is a system of integral rational numbers $c_{h,k}^{(r)}$ such that
\begin{equation}
 \alpha_h\beta_k=\sum_r c_{h,k}^{(r)}\gamma_r,\tag{3}
\end{equation}
and since conversely each number in $\mathfrak c$ arises by multiplying one number from $\mathfrak a$ and one from $\mathfrak b$ and adding these products (Vol. II, § 169), one also has
\begin{equation}
 \gamma_s=\sum_{h,k}a_{h,k}^{(s)}\alpha_h\beta_k,\tag{4}
\end{equation}
where the $a_{h,k}^{(s)}$ are likewise integral rational numbers.

Substituting (3) in (4) gives
\begin{equation}
 \gamma_s=\sum_r\gamma_r\sum_{h,k}a_{h,k}^{(s)}c_{h,k}^{(r)}\tag{5}
\end{equation}
and hence
\begin{equation}
 \sum_{h,k}a_{h,k}^{(s)}c_{h,k}^{(r)}=(r,s),\tag{6}
\end{equation}
where $(r,r)=1$, $(r,s)=0$, $r\ne s$.

Thus let $x_r$ denote variables and put
\begin{equation}
 \sum_r x_rc_{h,k}^{(r)}=y_{h,k}.\tag{7}
\end{equation}
It follows that
\begin{equation}
 x_s=\sum_{h,k}a_{h,k}^{(s)}y_{h,k}.\tag{8}
\end{equation}
Now let $u_r,v_r,t_r$ denote three systems of variables and
\begin{equation}
\begin{aligned}
 \mu&=\sum\alpha_ru_r,\\
 \nu&=\sum\beta_rv_r,\\
 \lambda&=\sum\gamma_rt_r
\end{aligned}\tag{9}
\end{equation}
basis forms of $\mathfrak a,\mathfrak b,\mathfrak c$, and let $U,V,T$ be the forms belonging to $\mathfrak a,\mathfrak b,\mathfrak c$, written in the variables $u,v,t$. Then
\begin{equation}
\begin{aligned}
 N(\mu)&=N(\mathfrak a)U,\\
 N(\nu)&=N(\mathfrak b)V,\\
 N(\lambda)&=N(\mathfrak c)T.
\end{aligned}\tag{10}
\end{equation}
From (3), by multiplication with $u_hv_k$ and summation according to (9), one obtains
\begin{equation}
 \mu\nu=\sum_r\gamma_r\sum_{h,k}c_{h,k}^{(r)}u_hv_k,\tag{11}
\end{equation}
so that, if
\begin{equation}
 t_r=\sum_{h,k}c_{h,k}^{(r)}u_hv_k\tag{12}
\end{equation}
is put,
\begin{equation}
 \mu\nu=\lambda,\tag{13}
\end{equation}
and since $N(\mathfrak a)N(\mathfrak b)=N(\mathfrak c)$, it follows by (10) that
\begin{equation}
 T=UV.\tag{14}
\end{equation}
Here, however, the $t_r$ are no longer independent variables; rather, they arise from $u_h,v_k$ by the bilinear substitution (12).

5. The form $T$ passes, by the bilinear substitution (12), into the product of the two forms $U,V$. This substitution, however, has one more essential peculiarity. We call the functions $t_r$ determined by (12) linearly independent modulo a prime number $p$ if from the congruence
\begin{equation}
 \sum x_rt_r=0\pmod p,\tag{15}
\end{equation}
where the $x_r$ are integral rational numbers, it follows that all these numbers are divisible by $p$. If this is the case for every arbitrary prime number $p$, the $t_r$ are simply called linearly independent.

The congruence (15), by (12) and (7), is equivalent to the $n^2$ congruences
\[
 y_{h,k}=\sum_r x_rc_{h,k}^{(r)}\equiv0\pmod p,
\]
and then (6) gives
\[
 x_s\equiv0\pmod p.
\]
This says that the substitutions (12) are linearly independent for every modulus $p$.

According to Gauss (Disq. aritm. art. 235), a form $T$ is called composed or compounded from the two forms $U,V$ if $T$ passes into the product $UV$ by a bilinear substitution for the variables $t$ whose equations are linearly independent for every modulus.

Consequently we have the theorem:

6. The form which belongs to the product of two ideals $\mathfrak a,\mathfrak b$ is composed from the forms of the ideals $\mathfrak a$ and $\mathfrak b$.
\footnote{The theorems which Gauss proves at the cited place for binary quadratic forms by a very extensive calculation are here derived in the greatest generality by simple considerations. Cf. Dirichlet-Dedekind, Lectures on Number Theory, § 182. For binary quadratic forms: Dedekind, Crelle's Journal, vol. 129; H. Weber, Göttinger Nachrichten 1907.}

\end{document}
